\section{Methods and modelled assumptions}
\label{app:methods}

This appendix records the ephemeris handling, frame and timescale conventions, the
sampling and numerical method, and the frame-tie discrimination behind the results.
All are consistent with the vendored provenance notice
(Section~\ref{sec:availability}).

\subsection{Ephemeris provenance and access}
\label{app:provenance}

Three independent authoritative ephemeris families are differenced, read through the
IMCCE \code{calceph} library so that all three are accessed by a single interface in
one frame and timescale:
\begin{itemize}
  \item \textbf{DE440} (JPL/NASA), kernel \code{de440s.bsp}
  \citep{park2021de440}, from
  \url{https://naif.jpl.nasa.gov/pub/naif/generic_kernels/spk/planets/de440s.bsp};
  \item \textbf{INPOP21a} (IMCCE / Observatoire de Paris), kernel
  \code{inpop21a\_TDB\_}\allowbreak\code{m100\_p100\_}\allowbreak\code{littleendian.dat}
  \citep{fienga2021inpop21a}, from
  \url{https://ftp.imcce.fr/pub/ephem/planets/inpop21a/};
  \item \textbf{EPM2021} (IAA~RAS), kernel \code{epm2021.bsp}
  \citep{pitjeva2021epm2021}, from \url{https://ftp.iaaras.ru/pub/epm/EPM2021/SPICE/}.
\end{itemize}
Following the vendoring convention of the companion paper's LLR fixtures
\citep{baweja2026datum}, the kernels are \emph{not} redistributed (they are large and
their redistribution terms vary); only the sampled derived positions and the computed
reference decomposition are vendored. Published scientific ephemeris positions are
factual constants and are not copyrightable; the derived sample points are cited with
full attribution to JPL, IMCCE, and IAA~RAS.

\subsection{Frame, timescale, and sampling}
\label{app:sampling}

All positions are computed in the ICRF/J2000 frame and Barycentric Dynamical Time
(TDB), so the three families are directly comparable without a frame or timescale
conversion between them. Two position sets are sampled over the window 2024-01-01 to
2025-12-31 (TDB) at two-day cadence ($366$ epochs):
\begin{itemize}
  \item \code{moon\_geo.csv}: the geocentric Moon position (NAIF body $301$ with
  respect to $399$), in metres. This is the object whose disagreement is decomposed
  in Table~\ref{tab:anchor}; its metre-per-nanoradian lever arm is the Earth--Moon
  distance $\Rmoon=\SI{3.840e8}{m}$.
  \item \code{planet\_ssb.csv}: the Mercury, Venus, Mars, and Earth--Moon-barycentre
  positions with respect to the Solar System barycentre (NAIF bodies $1,2,4,3$ with
  respect to $0$), in metres. These witness the common ICRF frame-tie rotation
  (Section~\ref{app:frametie}).
\end{itemize}
Both files, and the SciPy oracle, round all positions to six decimals (micrometre
resolution) \emph{before} the fit, so that the engine and the oracle consume
byte-identical inputs. The validation therefore tests agreement of the two fitting
\emph{algorithms}, not the coincidence of two independent samplings; any
sampling-level difference is removed by the shared rounded fixture.

\subsection{Numerical method}
\label{app:numerics}

The constant Helmert datum $\datum_{ab}$ of \eqref{eq:diffmodel} is the least-squares
solution of the stacked point-difference system with design matrix rows
$\Jac(\bm{p}_i)$. We solve the centred normal equations---subtracting the per-provider
centroid before assembling the system---with Jacobi (diagonal) preconditioning. Two
reasons motivate this choice. First, centring removes the translation--rotation
cross-coupling that otherwise inflates the condition number: for points at radius
$\sim\!\Rmoon$, the rotation and scale columns of $\Jac$ carry magnitudes
$\sim\!\Rmoon$ while the translation columns are $O(1)$, so the raw normal matrix has
condition number $\sim\!\Rmoon^2\sim10^{17}$; centring and Jacobi preconditioning
(dividing each column by its norm) bring the columns to comparable scale and the
condition number to $O(1)$, restoring double-precision accuracy. Second, for the
rotation-only sub-fit \eqref{eq:rotonly} we use the closed-form singular value
decomposition of the centred cross-covariance \citep{arun1987leastsquares}, which is
the numerically stable route to the rotation and is the form the independent oracle
also uses (so the two agree to the tight tolerance of Section~\ref{sec:oracle}). The
general-purpose linear algebra follows standard practice
\citep{golub2013matrix}.

\subsection{Frame-tie discrimination and the Sun exclusion}
\label{app:frametie}

The reducible/irreducible split requires separating a rotation common to all bodies
(a frame-tie) from a rotation specific to the Moon (dynamics). We estimate the
body-common rotation $\rot_{\mathrm{tie}}$ by fitting the same rotation-only model
\emph{separately} to each of the four witness bodies (Mercury, Venus, Mars, and the
Earth--Moon barycentre), obtaining four per-body rotation vectors, and taking
$\rot_{\mathrm{tie}}$ as their \emph{component-wise median} (the median of the four
$x$-components, and likewise for $y$ and $z$). This is a deliberately robust but
ad-hoc aggregate: the median resists a single discrepant witness, but it is
\emph{not} a joint fit across the witness bodies and \emph{not} a weighted fit, and a
different aggregation would shift the recovered $\rot_{\mathrm{tie}}$---and hence the
split---quantitatively (Section~\ref{sec:limitations}). The Moon-specific excess is
then $\rot_{\mathrm{exc}}=\rot_{\mathrm{moon}}-\rot_{\mathrm{tie}}$, and the metre
envelopes of Table~\ref{tab:anchor} follow by mapping $\rot_{\mathrm{tie}}$ and
$\rot_{\mathrm{exc}}$ through the Earth--Moon lever arm; because the two are in general
non-parallel, these are worst-case upper bounds that do not sum to the raw
disagreement.

The witness set is the four inner bodies Mercury, Venus, Mars, and the Earth--Moon
barycentre. The Sun is deliberately \emph{excluded}. The Sun's position with respect
to the Solar System barycentre is dominated by the reflex motion imposed by the giant
planets and is, by definition, tied to the barycentre convention itself; a difference
in how two ephemeris families realize the barycentre therefore contaminates the Sun's
SSB position with a signal that is not a clean rigid frame rotation. Including it would
bias the frame-tie estimate toward a barycentre-definition artifact rather than the
ICRF orientation tie we intend to capture. The inner planets and the Earth--Moon
barycentre, by contrast, have SSB positions dominated by their own orbits and provide
cleaner witnesses of a common frame rotation. This choice is a \Modelled{} decision;
Section~\ref{sec:limitations} notes that a different witness set or weighting would
shift the split quantitatively while leaving the qualitative Moon-excess dominance
(for the JPL-versus-others pairs) intact, and that the convention-free content is the
Moon-excess itself (Remark~\ref{rem:attribution}).
